We present \textsc{bisect-apportion}, a Rust crate that currently implements
four divisor-method apportionment paths --- Huntington-Hill, Webster, Adams,
and Jefferson/d'Hondt --- with two primary public functions:
\texttt{huntington\_hill()} and \texttt{apportionment\_divisor()}.
Hamilton/largest-remainder apportionment is a paper and paradox-analysis
reference point, not a shipped public apportionment function in this crate.
The central implementation claim is narrower than earlier drafts: the crate
exposes deterministic priority-queue divisor methods with constitutional-floor
handling, deterministic tie-breaking, and executable unit/property tests.

The priority computations use \texttt{f64} arithmetic.
For the seat counts and population values in U.S.\ apportionment
(populations $\leq 40$~million, seat counts $\leq 100$), the observed
priority gaps are far larger than floating-point rounding error, so the
current fixture allocations remain stable.  Full Census Bureau regression
fixtures and SHA-256 publication verification remain future hardening steps;
the current crate does not yet ship a cryptographic apportionment verifier or
a \texttt{bisect apportion} CLI command.

The test suite comprises 99 tests across the crate:
L0 (unit tests on priority formulas and edge cases),
L1 (integration tests on synthetic examples with known correct answers),
and property tests for prime-factor, nesting, spectral, and split helpers that
share the crate.  Twenty-four tests specifically cover apportionment methods
and paradox checking.
The crate also implements \texttt{check\_alabama\_paradox()},
which empirically checks house monotonicity over caller-provided house-size
sequences for divisor methods.

\textsc{bisect-apportion} is the apportionment and prime-factor compositor
component of the \textsc{bisect} redistricting platform, providing reusable
apportionment functions and redistricting tree helpers used elsewhere in the
platform.
